%% 第 5 章  自动生成于 scripts/tex/md2tex.py，请勿直接手改
\chapter{基于随机化辅助的 ARM 指针认证机制}
\label{chap:5}

\section{引言}

研究内容二给出了一种低开销的触发式随机化方法，但它同时为内核地址空间留下一片非随机化区域，其中的布局在运行期不发生变化。这并非某一种方法的缺陷：受开销所限，现有的内核持续随机化方法全部只能做部分随机化，Adelie\citep{nikolaev2022adelie}、Chaos\citep{chang2024chaos} 与 Dbox\citep{li2024dbox} 以模块与驱动为界，随机化区域之外的核心内核保持静态。非随机化部分的防护因此是所有这类方法都必须面对的问题。本章旨在给出一种与具体随机化方法解耦的非随机化区域保护方法。

布局长期不变意味着攻击者一旦持有一个合法的代码指针值，它就长期可用。

理想的防御手段是硬件辅助的控制流完整性，其中 ARM 指针认证的上下文由软件任意选取，原理上可以为每个调用点取不同的值，但受上下文传播约束：多个调用点调用同一函数时必须共享上下文，该约束沿调用图传播，最终使大量调用点收敛到同一上下文。本章的关键观察是：随机化恰好提供了打断这条传播链的手段。随机化区域内的目标地址不可预测，进入该区域的间接调用不需要指针认证；把造成传播的高扇入函数移入随机化区域，链两侧的调用点即可各自建立唯一上下文。

本章的组织如下。第 5.2 节说明部分随机化留下的非随机化区域、指针认证的上下文传播问题，并归纳挑战；第 5.3 节给出系统设计，包括总体架构与核心思想、以随机化打断上下文传播链的识别与绑定模型、以及跳板处的分区验证；第 5.4 节按编译期、运行期与出口检查三部分给出实现；第 5.5 节评估；第 5.6 节给出安全性分析；第 5.7 节小结。

\section{研究背景与动机}

\subsection{部分随机化与非随机化区域}

现有的内核持续随机化方法以模块与驱动为随机化对象，核心内核保持静态。Adelie 对非随机化区域的处理是把指向它的指针约束为只能经只读全局偏移表访问，指向模块内部的指针则在模块移动后自动失效。这一做法依赖模块与内核之间有限且明确的接口，在 built-in 场景下没有一张表能把跨区域引用收拢。

更一般的处理属于控制流完整性的范畴。软件实现的控制流完整性在内核中的开销不可接受；硬件辅助的实现中，落地点标记类机制（Intel CET-IBT、ARM BTI）只检查目标处有无标记\citep{shanbhogue2019cet,arm2023armv8}，把所有被标记的函数放进同一个等价类；签名验证类机制（ARM 指针认证）原理上可以做到每个调用点一个上下文，其实际精度取决于上下文如何选取。

\subsection{指针认证与上下文传播}

ARM 指针认证在指针的未使用高位中嵌入密码学签名，签名的输入包括指针值、一个 64 位上下文与一个软件不可读的硬件密钥。签名与验证由专用指令完成，各约四个周期，认证码随指针搬移而自然传播，不需要软件协助。图~\ref{fig:5-1} 给出这一对运算的全过程：指针赋值处以密钥与上下文对指针值计算认证码并写入高位，指针使用处以同一密钥与同一上下文重算并比对，一致则还原出可用的指针，不一致则高位被置为错误码，随后的解引用触发异常。密钥由硬件保存且软件无法读出，攻击者因此无法离线伪造认证码；上下文则完全由软件给出，它决定一个指针能够在哪些位置通过认证。

\begin{figure}[!htbp]
    \centering
    \includegraphics[width=\textwidth,height=0.46\textheight,keepaspectratio]{schematic/ch5-pac-primitive}
    \bicaption{\enspace 指针认证的签名与认证过程}{\enspace Signing and Authentication in ARM Pointer Authentication}
    \label{fig:5-1}
\end{figure}

指针认证的精度完全取决于上下文的选取。若每个调用点使用互不相同的上下文，则在某调用点签名的指针无法在另一调用点通过验证，等价类大小为一。这一理想状态无法达到，原因是上下文传播：当多个调用点可能调用同一函数时，这些调用点必须使用相同的上下文，否则被调用方无法验证；该约束沿调用图向外扩散，最终把大量调用点并入同一上下文。图~\ref{fig:5-2} 给出这一过程：两个调用点共用一个被调函数，就必须共享上下文，该约束沿调用图扩散，把越来越多的调用点并入同一集合。

\begin{figure}[!htbp]
    \centering
    \includegraphics[width=\textwidth,height=0.46\textheight,keepaspectratio]{schematic/ch5-context-propagation}
    \bicaption{\enspace 指针认证的上下文传播问题}{\enspace The Context-Propagation Problem of Pointer Authentication}
    \label{fig:5-2}
\end{figure}

PARTS\citep{liljestrand2019pac}、PACStack\citep{liljestrand2021pacstack}、PACTight\citep{ismail2022pactight} 等工作分别改进上下文的构造方式或扩展保护对象，面向内核的指针认证应用亦有若干工作。这些工作都未能从根本上打断传播链，因为传播的根源是调用图的结构本身：只要一个函数同时被两个调用点调用，且它需要验证传入的指针，两个调用点就必须共享上下文。在纯指针认证的框架内，这个约束难以解除。本文在同一平台上按函数类型哈希取上下文的方案，最大等价类为 141（第 5.5.2 节）。

\subsection{研究动机与挑战}

上两节归纳出两点：非随机化区域的保护是所有部分随机化方法的共同问题；商用硬件控制流完整性机制中，指针认证是唯一能以软件选取的上下文区分调用点、因而精度上限最高的一种，但它受上下文传播所限而实际退化。打断传播链需要引入指针认证之外的信息，本章引入的是随机化。

本章面对的挑战有三。其一，如何在不做全局指针分析的前提下让调用点得到唯一上下文，内核规模决定了任何依赖全局点到分析的方案都不具备可部署性。其二，如何避免在每一次间接调用上都施加验证，使开销随随机化范围可控。其三，如何使保护机制不绑定某一种随机化实现，从而可以施加到任何一种部分随机化方法留下的非随机化区域上。

\section{系统设计}

\subsection{总体架构与核心思想}

本章的设计可以概括为：进入随机化区域的控制流以范围检查替代认证，因为目标地址不可预测；造成上下文传播的高扇入函数移入随机化区域，传播链在此剪断；只对离开随机化区域的控制流施加认证与白名单的双重检查。验证施加在跨区域转移的唯一出入口上，不另建检查点。本文原型中这一出入口即第 3 章的跳板。图~\ref{fig:5-3} 按时间顺序给出机制的三个阶段。启动早期为识别阶段，指针以类型上下文签名，归属记录跟踪每个函数的首个调用点，第二个调用点触及即判为高扇入。启动推进中为绑定与迁移阶段，单一归属的函数重签为调用点唯一上下文，高扇入函数移入随机化区域，指向它的指针改写为跳板地址。运行期为分区验证阶段，入口只验范围，出口验认证与白名单。绑定后新出现的归属冲突经仲裁后同样迁移，图中底部的回边表示这一情况。

\begin{figure}[!htbp]
    \centering
    \includegraphics[width=\textwidth,height=0.46\textheight,keepaspectratio]{schematic/ch5-architecture-mono}
    \bicaption{\enspace 本章机制的总体架构}{\enspace Overall Architecture of the Proposed Mechanism}
    \label{fig:5-3}
\end{figure}

\subsection{以随机化打断上下文传播链}

\textbf{识别。} 识别高扇入函数无法被静态分析，它可以作为验证流程的自然副产品被观测到。

本章采用指针在存储中始终以类型上下文签名，另以一条归属记录跟踪首个调用它的调用点。类型签名是合法性的真值来源，存储中的函数指针恒以函数原型的哈希为上下文签名，任何调用点都能以类型上下文认证其合法性，认证失败即为篡改。归属记录识别扇入：首个使用该指针的调用点将其认领，同一调用点再次使用属正常，另一个调用点使用同一指针即判定为高扇入。识别只读归属记录，不改写指针本身的签名，因而不会陷入上述矛盾。

\textbf{上下文绑定。} 精度来自上下文绑定，识别只负责分类。上述识别把函数分成单一归属与高扇入两类。对单一归属的函数，在识别稳定后由绑定阶段把它改签为其归属调用点的唯一上下文并不再变更。绑定之后等价类以调用点为单位划分：每个函数至多归属一个调用点，一个调用点却可以归属多个函数，它经同一指针变量先后调用的全部函数都以它的上下文签名，其合法目标集合就是这些函数。单目标调用点的等价类为一，分发类调用点的等价类等于它分发到的函数个数，这是调用图的固有结构而非机制的损失。对高扇入函数则移入随机化区域，此后不再参与验证，传播链在此处被切断。绑定之后调用点的等价类不再增长：绑定后新出现的目标，无论是新的调用点使用已绑定的函数，还是已绑定的调用点使用新的函数，一律交由随机化承担，处理路径见第 5.4.2 节。

一个函数一旦被第二个调用点调用，它就不再需要指针认证，改由持续随机化保护，这条规则对绑定前后同样成立。绑定前的第二个调用点由归属记录直接判出。绑定后的第二个调用点无法通过认证，因为内存中的指针只携带第一个调用点的认证码，它既通不过自己的上下文也通不过类型上下文，认证本身无法区分“绑定后新出现的合法调用点”与“指针被篡改”。本章把这一区分交给归属记录：认证失败进入违规处理路径，处理路径按指针的存储位置查记录，若记录已绑定、且指针的值等于记录中的目标函数，则判定为绑定后的归属冲突，该函数即刻判为高扇入并进入随机化；指针与记录不符则维持违规判定。转换完成后，指向该函数的指针被改写为其入口跳板的地址，不再携带认证码，此后无论第一个还是第二个调用点，调用它都只做范围检查。图~\ref{fig:5-4} 给出一次识别的时序、绑定的时间点与绑定后归属冲突的处理：指针签名在识别阶段恒为类型上下文，归属由独立记录跟踪，归属冲突即判为高扇入并移入随机化区域，单一归属的函数在绑定时重签为调用点唯一上下文。识别贯穿内核启动过程：启动早期所有指针以类型上下文签名以保证兼容，随着执行推进、归属趋于稳定，单一归属的函数被逐步绑定，启动完成后大部分单一归属的函数已建立唯一上下文。

\begin{figure}[!htbp]
    \centering
    \includegraphics[width=\textwidth,height=0.46\textheight,keepaspectratio]{schematic/ch5-fanin-identification}
    \bicaption{\enspace 高扇入函数的识别与绑定}{\enspace Identification and Binding of High-Fan-In Functions}
    \label{fig:5-4}
\end{figure}

\textbf{信任根。} 唯一上下文的建立依赖启动阶段不存在控制流劫持：若攻击者在启动阶段即可劫持控制流，就可以构造恶意调用点，用该调用点的上下文为任意目标签名，此后即可用该上下文通过验证。该假设由安全启动链保证，绑定因此落在第一个用户态进程启动之前。此外假设攻击者在劫持成功之前不能任意改写内存权限，使随机化过程中更新的代码页不会被预先篡改。

\subsection{跳板处的分区验证}

验证按控制流的方向差异化处理。

\textbf{进入方向以范围检查替代认证。} 随机化区域内的目标地址不可预测，攻击者无法把指针改写为区域内某个函数体的地址。但对函数取地址得到的是跳板地址，而跳板位置固定且公开，攻击者把函数指针改写为某个入口跳板即可确定性地到达对应函数的入口，因此入口不能完全免除检查。本章在每个间接调用点先判断目标是否落在跳板区域之内：落在其中，说明目标函数处于随机化状态，只做这一次范围检查；不在其中，目标位于非随机化区域，以调用点上下文认证。跳板区域位置固定、边界在链接期确定，这是一次常数时间的边界比较，代价远低于一次认证指令加上下文计算。指针的形态由此决定检查方式：带认证码的指针走认证，跳板地址走范围检查。一个函数由非随机化状态转为随机化状态时，只需把指向它的指针改写为跳板地址，不需要改动任何调用点的指令。它保证的性质是以调用形式进入随机化区域的控制流只能落在函数入口，因为跳板是调用进入该区域的唯一形式，而每个跳板对应一个函数的入口。返回方向经返回标记落到函数内偏移，其落点在函数中部而非入口，且不经范围检查；这条路径的约束来自标记本身的完整性，标记以栈槽地址为上下文签名，搬到别的栈槽即认证失败。

\textbf{离开方向强制认证加白名单。} 目标位于非随机化区域，地址长期不变。出口依次执行：以调用点唯一上下文认证目标指针，检查目标属于编译期确定的合法外部函数集合，两项都通过才放行。两项检查约束的对象不同：白名单约束目标是否属于合法集合，属于集合成员性检查；指针认证约束这个指针是否由本调用点签发，属于来源真实性检查。仅有白名单时攻击者可以把一个合法目标替换为白名单内的另一个合法目标，仅有指针认证时攻击者可以复用一个在别处合法签发的指针，两者叠加才构成完整约束。图~\ref{fig:5-5} 对照两个方向的处理：入口只验范围，保证落点为函数入口；出口验认证与白名单，一验来源一验集合。

\begin{figure}[!htbp]
    \centering
    \includegraphics[width=\textwidth,height=0.46\textheight,keepaspectratio]{schematic/ch5-partitioned-check}
    \bicaption{\enspace 跳板处的分区验证}{\enspace Partitioned Verification at the Trampolines}
    \label{fig:5-5}
\end{figure}

\textbf{对所承接随机化原语的要求。} 免验的前提是目标地址不可预测，因此免验范围严格等于当前处于随机化状态的函数集合，而非编译期划定的区域。本章据此对随机化原语只提出一项要求：被识别为高扇入的函数必须持续处于随机化状态。本文原型中，第 4 章的激活集合恒含本章的高扇入集合。这一要求之外，本章不依赖随机化原语的任何内部细节，它可以施加在任何提供函数粒度随机化区域与跨区域出入口的方案之上。

\section{系统实现}

\subsection{编译期：指针认证插桩}

本章在 LLVM Pass 上扩展两类插桩：函数指针产生处插入签名指令，初始使用类型上下文；间接调用处先做一次跳板区域的范围检查，目标落在其中即直接调用；否则插入认证指令，优先使用调用点唯一上下文，失败则回退到类型上下文。类型上下文的回退在绑定之前覆盖识别阶段，在绑定之后则说明该指针是绑定后新产生的；回退也失败则说明指针携带的是另一个调用点的认证码或已被篡改。两种失败都进入违规处理路径，由第 5.4.2 节的归属仲裁决定放行、原位重签、纳入随机化或判定违规。

上下文是指针认证指令的修饰值，它决定一个签名指针只能在哪些位置通过验证，因而是 CFI 精度的载体。上下文在使用处现算，不占用任何长期驻留的寄存器：类型上下文取函数原型的哈希，用于识别阶段中所有存储于内存的指针；调用点上下文取调用指令虚拟地址的哈希，用于绑定之后单一归属函数的签名与验证。两者都是对一个已知输入做一次常数时间的哈希，结果直接作为 \lstinline[breaklines=false]@PACIA@ 与 \lstinline[breaklines=false]@AUTIA@ 的修饰值，签名与验证不需要读取任何全局状态。

构造上有两处细节。其一是域分离：类型上下文的哈希输入中置一个固定的区分位，使两类上下文的输入空间不相交，否则一个以类型上下文签名的指针有可能恰好通过某个调用点上下文的验证。其二是非零约束：两类上下文都保证取值非零，避免修饰值为零这一退化情形与上下文尚未设置的状态相混。不把上下文基址长期驻留在寄存器中有两重理由：被调用者保存的寄存器长期占用会与内核中的手写汇编及既有调用约定冲突；驻留寄存器的基址一旦被读到即可复用。上下文不驻留于任何位置，攻击者即便读取全部通用寄存器也得不到一个可复用的基址，伪造一个能通过验证的指针仍然需要密钥本身。

\subsection{运行期：高扇入函数的跟踪、绑定与迁移}

跟踪以每个函数指针一条记录的形式组织，记录中保存指针所在的存储位置、目标函数、其类型哈希、归属调用点、以及两个状态位：是否已判定为高扇入、是否已绑定。某个函数指针首次被某调用点使用且通过类型上下文验证时，该调用点被登记为它的归属；此后同一调用点再次使用属于正常；一旦另一个调用点使用同一指针，归属冲突，该函数即被判定为高扇入。

绑定之后，调用点唯一上下文认证失败的指针分两种情形，都在违规处理路径上按归属记录仲裁，算法~\ref{alg:5-1} 给出间接调用处从范围检查到仲裁的完整判定过程。其一，类型上下文回退也失败，指针携带的是另一个调用点的认证码，或已被篡改。处理路径按指针的存储位置查记录：记录不存在或未绑定，即判定为控制流完整性违规并终止；记录已绑定且函数指针等于目标函数 F，则判定为绑定后的归属冲突，本次调用以 F 为目标直接放行，同时把 F 排入迁移队列。其二，类型上下文回退通过，指针是绑定之后新产生的，携带类型认证码。若记录已绑定、归属为本调用点、且函数指针等于记录中的目标函数，这只是同一指针值被重新写入，处理路径在原位把它重签为本调用点的唯一上下文，等价类不变；否则，无论是本调用点指向了新的函数，还是另一个调用点使用了该指针，目标函数都排入迁移队列。类型认证通过是纳入随机化的前提，它证明该指针由合法的签名插桩产生，目标是编译器允许出现在这一类型槽位上的函数；否则攻击者只需写入一个裸指针，就能把任意函数送入随机化区域并经跳板调用它。

两种情形下处理路径都不分配内存、不睡眠，可在任意上下文执行。随机化线程在关键路径之外完成排队函数的物理迁移，并把归属记录中登记的、指向它的全部指针存储位置改写为其入口跳板地址，不带认证码，此后所有调用点对这些指针都只做范围检查。在迁移完成之前，每次触发的调用都要经处理路径放行，这是一段微秒级的过渡。取函数地址的一侧也随之一致：迁移后对该函数取地址得到的本就是跳板地址，产生指针处的签名插桩对落在跳板区域内的地址不签名，新产生的指针天然是范围检查的形态。绑定之后调用点的等价类由此只减不增，新目标一律由随机化承担，而不并入任何调用点的上下文。

\begin{algorithm}[!htbp]
    \caption{间接调用处的检查与归属仲裁}
    \label{alg:5-1}
    \begin{algorithmic}[1]
        \Require 调用点 $CS$，函数指针的存储位置 $loc$，其中的指针值 $p$
        \Ensure 调用目标 $t$，或判定为控制流完整性违规
        \Function{CheckAndArbitrate}{$CS, loc, p$}
            \If{$p \in [\mathit{tramp}_{lo}, \mathit{tramp}_{hi}]$} \Comment{指向随机化区域：只做范围检查}
                \State \Return $p$
            \EndIf
            \State $t \gets \Call{Autia}{p, H(\&CS)}$ \Comment{调用点唯一上下文}
            \If{$t$ 有效} \State \Return $t$ \EndIf
            \State $t \gets \Call{Autia}{p, H(\mathit{type})}$ \Comment{回退到类型上下文}
            \State $r \gets \mathit{Record}[loc]$ \Comment{归属仲裁：按存储位置查记录}
            \If{$t$ 有效} \Comment{指针由合法的签名插桩产生}
                \If{$r$ 未绑定} \Comment{识别阶段}
                    \If{$r.\mathit{owner} = \varnothing$} $r.\mathit{owner} \gets CS$
                    \ElsIf{$r.\mathit{owner} \neq CS$} \Call{Migrate}{$t$} \Comment{第二个调用点：高扇入}
                    \EndIf
                    \State \Return $t$
                \EndIf
                \If{$r.\mathit{owner} = CS$ \textbf{and} $t = r.\mathit{target}$} \Comment{同一指针值重新写入}
                    \State $\mathit{Mem}[loc] \gets \Call{Pacia}{t, H(\&CS)}$ \Comment{原位重签，等价类不变}
                    \State \Return $t$
                \EndIf
                \State \Call{Migrate}{$t$}; \Return $t$ \Comment{绑定后的新目标：纳入随机化}
            \Else \Comment{携带其他调用点的认证码，或已被篡改}
                \State $t \gets \Call{Strip}{p}$
                \If{$r$ 已绑定 \textbf{and} $t = r.\mathit{target}$} \Comment{绑定后的归属冲突}
                    \State \Call{Migrate}{$t$}; \Return $t$
                \EndIf
                \State \Return 违规
            \EndIf
        \EndFunction
        \Statex
        \Procedure{Migrate}{$F$} \Comment{不分配内存、不睡眠，可在任意上下文执行}
            \State 标记 $F$ 为高扇入，排入迁移队列
            \State 随机化线程在关键路径之外完成物理迁移，并把 $\mathit{Record}$ 中指向 $F$ 的全部存储位置改写为 $F$ 的入口跳板地址
        \EndProcedure
    \end{algorithmic}
\end{algorithm}

归属记录在这条路径上成为合法性的仲裁依据，因此它的完整性是前提：记录在绑定完成后置为只读，只在迁移路径上临时开放写入；攻击者若把指针改写为另一个函数 G，函数指针与记录不符，仍被判为违规。

\subsection{跨区域出口的双重检查}

白名单约束从随机化区域到非随机化区域的合法目标。表项来自链接器汇总的专用数据段，编译器扫描随机化代码的目标文件，提取以外部符号或未定义符号形式出现的跨边界控制流目标；内核初始化时把它们装入哈希表，随后将该表置为只读。哈希表的桶数取不小于表项数两倍的、最接近的 2 的幂，使平均链长小于一，查询为哈希定桶加桶内线性比对，因而是常数时间。在白名单尚未完成装入之前查询一律返回通过，以免阻断早期启动路径；这段窗口位于用户态启动之前，与第 5.3.2 节所依赖的安全启动链假设处在同一时段之内。

出口先以调用点唯一上下文认证目标指针，再查询白名单，两者都通过才放行。前者检查指针是否由本调用点签发，后者检查目标是否属于合法集合，二者不可互相替代。

\section{实验评估}

本节回答两个问题：随机化打断上下文传播链之后，控制流完整性检查的精度提高到什么程度、代价多大；这一机制是否真的与具体的随机化方法解耦，会对现有随机化方法产生什么影响。

\subsection{实验配置}

ARM Neoverse N2，FEAT\_PAuth，Linux 6.8，LLVM 14。等价类统计在 x86-64 上按目标文件的重定位与调试信息完成，与平台无关。运行期快照以 UnixBench 一次完整运行驱动，开销以 LMBench 与 SPEC CPU2006 测量。10 次运行取中位。

\subsection{控制流完整性检查的精度与开销}

精度以每个间接调用点的合法目标数衡量：在给定机制下，一个间接控制流转移有多少个目标可选。表~\ref{tab:5-1} 汇总已发表的内核控制流完整性机制与本文的结果，末三行为本文在 Linux 6.8 上的统计，其余各行取自所引原文；其中 BTI 与 Clang CFI 两行的数值由 PAL 在其表 2 中一并给出。

\begin{table}[!htbp]
    \bicaption{\enspace 合法目标数分布对比}{\enspace Comparison of Distribution of Legitimate Targets per Indirect Call Site}
    \label{tab:5-1}
    \centering
    \scriptsize
    \setlength{\tabcolsep}{3pt}
    \renewcommand{\arraystretch}{1.25}
    \hyphenpenalty=10000\exhyphenpenalty=10000\emergencystretch=1em
    \begin{tabular}{>{\raggedright\arraybackslash}p{\dimexpr 0.2700\textwidth-2\tabcolsep\relax}>{\raggedright\arraybackslash}p{\dimexpr 0.1800\textwidth-2\tabcolsep\relax}>{\raggedleft\arraybackslash}p{\dimexpr 0.1350\textwidth-2\tabcolsep\relax}>{\raggedleft\arraybackslash}p{\dimexpr 0.1350\textwidth-2\tabcolsep\relax}>{\raggedleft\arraybackslash}p{\dimexpr 0.1250\textwidth-2\tabcolsep\relax}>{\raggedleft\arraybackslash}p{\dimexpr 0.1550\textwidth-2\tabcolsep\relax}}
        \hline
        机制 & 内核 & $\leqslant 5$ & $\geqslant 100$  & 最大 & 平均目标数\\
        \hline
        BTI\citep{yoo2022pal} & Linux 5.12 & 0.0\% & 100\% & 59,300 & -\\
        Clang CFI\citep{tolvanen2018androidcfi} & Android 4.14 & 55.0\% & 7.0\% & 1,153 & 59\\
        PAL\citep{yoo2022pal} & Linux 5.12 & 84.9\% & 2.8\% & 35,264 & -\\
        HEK-CFI\citep{maar2024hekcfi} & Linux 5.18 & 65.3\% & 7.4\% & - & 22\\
        类型上下文 & Linux 6.8 & 55.0\% & 2.0\% & 141 & 13.7\\
        唯一上下文，启动完成 & Linux 6.8 & 91.0\% & 0\% & 57 & 3.6\\
        唯一上下文，运行期 & Linux 6.8 & 91.95\% & 0\% & 57 & 3.4\\
        \hline
    \end{tabular}
\end{table}

不超过 5 的占比在 Android 4.14 与本文的 Linux 6.8 上都是 55\%，而最大值相差一个数量级，1,153 与 141 分别来自 Google 的类型匹配和本文按地址被取过的函数收紧后的分组。HEK-CFI 只报告平均目标数，它把类型哈希在同一内核上的 325 降到 22，是近年内核控制流完整性中精度最高的一支，但其分布未见诸原文。这一组也说明本文选择类型上下文作为对照的理由：它是主线 kCFI\citep{clang2022kcfi} 与 FineIBT\citep{gaidis2023fineibt} 的划分依据，也是 PAL 未加精化前的起点。上下文精化一组把 PAL 的两级精化与本文的唯一上下文并列。PAL 通过静态分析把上下文从类型细化到对象类型，再对仍然过宽的调用点绑定对象实例与调用上下文，由分析器自动加入 166 处注释后，不超过 5 的占比升到 90.8\%，超过 100 的降到 0.08\%，仍有 3 个上下文超过 100 个目标，最大 207。本文以调用点地址为上下文，启动完成时不超过 5 的占比为 91.0\%，与 PAL 相差在 0.2\%以内，最大值为 57 且没有任何调用点超过 100，代价是把 233 个高扇入函数移入随机化区域。两者的差别在于高扇入调用点的处置：PAL 靠更细的静态上下文继续切分，切不开的留在大等价类里；本文不再切分，把这些函数交给随机化，等价类的残余上界由此从 207 降到 57。两种做法并不排斥，PAL 的静态精化可以作为本文识别阶段的初始上下文，本文机制只要求初始上下文可比较，不要求它是类型哈希。

其它以精度为主张的控制流完整性工作使用不同的指标，无法与上表并列。Ge 等人\citep{ge2016finegrained}在 FreeBSD 与 MINIX 内核上以 AIA 即每个间接转移允许目标数的平均值报告，kCFI 白皮书（Moreira 等人\citep{moreira2017droptherop}）同时报告 AIR 与 AIA；平均值把长尾折进均值里，Ge 等人自己也指出 AIR 无法反映过近似的程度，这正是本文同时报告分布、最大值与平均值的原因。OS-CFI\citep{khandaker2019oscfi}、HPCFI\citep{kasten2024hpcfi}、Burow 等人\citep{burow2017cfi}的综述与 CScan\citep{li2020cracks} 都以每个调用点的目标集合大小或其运行期可达数为指标，与本文口径同源，但保护对象是用户态程序，数值随程序而异，不能与内核数字并列。

本文类型上下文的最大值 141 见按表~\ref{tab:5-2}：最通用的几个函数原型聚集了较多函数，一个类型为 \lstinline[breaklines=false]@void (*)(void)@ 的间接调用在类型哈希机制下有 141 个合法目标可跳。类型分布的长尾性是以类型划分等价类这一选择的直接后果。本文按调试信息中的函数原型分组，把所有枚举折叠为一类、把匿名结构体折叠为一类、并去除类型别名，这比类型哈希机制实际的分组更粗，因此 141 是一个上界，该机制真实的最大等价类只会更小或相等。

\begin{table}[!htbp]
    \bicaption{\enspace 类型哈希机制下最大的几个等价类}{\enspace The Largest Equivalence Classes under the Type-Hash Mechanism (Tightened Criterion)}
    \label{tab:5-2}
    \centering
    \footnotesize
    \setlength{\tabcolsep}{3pt}
    \renewcommand{\arraystretch}{1.25}
    \begin{tabular}{lr}
        \hline
        类型签名 & 等价类大小\\
        \hline
        \lstinline[breaklines=false]@void (*)(void)@ & 141\\
        \lstinline[breaklines=false]@unsigned long long (*)(unsigned long long, ...)@ & 133\\
        \lstinline[breaklines=false]@int (*)(void)@ & 61\\
        \lstinline[breaklines=false]@int (*)(struct device*)@ & 57\\
        \lstinline[breaklines=false]@void (*)(struct sock*)@ & 55\\
        \hline
    \end{tabular}
\end{table}

本文机制的等价类以调用点为单位：调用点上下文由调用指令地址导出，互不相同，绑定之后每个调用点只接受由它自己签发的指针，其合法目标集合就是它所归属的全部函数。因此单目标调用点的等价类为一，而经同一指针变量分发到多个函数的调用点，等价类等于其目标函数的个数。被判定为高扇入的函数退出这一保护而由随机化承担，其规模由第 5.5.3 节给出。全内核规模的等价类分布按每个间接调用点的合法目标数统计，对照三档上下文策略：常数零上下文（等价于落地点标记）、按函数类型的上下文、本章唯一上下文。

统计取三个快照时刻。早期初始化指内核启动到第一个文件系统挂载之前，此时绑定尚未开始，全部指针仍以类型上下文签名。启动完成指第一个用户态进程启动之前绑定例程结束的时刻，它同时是第 5.3.2 节所依赖的信任根边界，绑定必须在此之前完成。运行期指启动完成后再完整执行一轮 UnixBench，其全部子项结束的时刻；选择这一时刻是因为它覆盖了进程创建、文件系统、管道与信号等启动期不会走到的冷路径，用以检验启动期得到的归属是否会被长期运行推翻。分布按目标数逐格统计。

图~\ref{fig:5-6} 以五条曲线给出各机制的分布：Clang CFI 的类型哈希、HEK-CFI\citep{maar2024hekcfi}、PAL 的类型签名，以及本章机制在启动完成与运行期两个快照下的结果。横轴为每个间接调用点的合法目标数，纵轴为取该目标数的调用点占比，末格汇总目标数超过 30 的调用点。各机制的主峰都落在 30 以内，横轴因此只画到 30，峰后的长尾按占比递减。峰后的次级隆起对应按文件系统类型、协议族或设备类别分发的调用点。末格的差距同样明显，目标数超过 30 的调用点在类型哈希下占 13\%，在本章机制下降到 2.2\%。落地点标记一类机制只有一个等价类，其分布退化为末格上的单点，不列入该图。HEK-CFI 以 Intel CET 影子栈的写保护页保存控制数据，是近年内核控制流完整性中精度最高的一支，每个间接调用的平均合法目标数为 22。表~\ref{tab:5-3} 给出三个快照的汇总。

\begin{figure}[!htbp]
    \centering
    \includegraphics[width=0.95\textwidth]{plots/ch5-cfi-precision-dist}
    \bicaption{\enspace 每个间接调用点合法目标数的分布}{\enspace Distribution of Legitimate Targets per Indirect Call Site}
    \label{fig:5-6}
\end{figure}

\begin{table}[!htbp]
    \bicaption{\enspace 三档上下文策略在三个快照时刻的精度}{\enspace Precision of the Three Context Policies at Three Snapshots}
    \label{tab:5-3}
    \centering
    \scriptsize
    \setlength{\tabcolsep}{3pt}
    \renewcommand{\arraystretch}{1.25}
    \hyphenpenalty=10000\exhyphenpenalty=10000\emergencystretch=1em
    \begin{tabular}{>{\raggedright\arraybackslash}p{\dimexpr 0.1086\textwidth-2\tabcolsep\relax}>{\raggedright\arraybackslash}p{\dimexpr 0.0931\textwidth-2\tabcolsep\relax}>{\raggedleft\arraybackslash}p{\dimexpr 0.2249\textwidth-2\tabcolsep\relax}>{\raggedleft\arraybackslash}p{\dimexpr 0.2172\textwidth-2\tabcolsep\relax}>{\raggedleft\arraybackslash}p{\dimexpr 0.0931\textwidth-2\tabcolsep\relax}>{\raggedleft\arraybackslash}p{\dimexpr 0.2482\textwidth-2\tabcolsep\relax}}
        \hline
        策略 & 快照 & 目标数不超过 5 的调用点占比 & 目标数超过 50 的调用点占比 & 最大等价类 & 移入随机化区域的函数数（占比）\\
        \hline
        常数零上下文 & 任意 & 0\% & 100\% & 44,777 & 不适用\\
        类型上下文 & 任意 & 55.0\% & 10.2\% & 141 & 不适用\\
        唯一上下文 & 早期初始化 & 55.0\% & 10.2\% & 141 & 0（0\%）\\
        唯一上下文 & 启动完成 & 91.0\% & 0.76\% & 57 & 233（4.4\%）\\
        唯一上下文 & 运行期 & 91.95\% & 0.66\% & 57 & 244（4.8\%）\\
        \hline
    \end{tabular}
\end{table}

表~\ref{tab:5-3} 与图~\ref{fig:5-6} 结合来看。早期初始化时全部指针仍以类型上下文签名，唯一上下文一行与类型上下文一行完全相同，55\% 的调用点目标数不超过 5，10.2\% 超过 50，最大 141。绑定随启动推进逐步完成，启动完成时不超过 5 的调用点占比升到 91.0\%，超过 50 的降到 0.76\%，最大等价类由 141 降到 57；这一变化只需把 233 个高扇入函数移入随机化区域，占间接调用目标函数的 4.4\%。运行期一次完整的 UnixBench 之后再识别出 11 个，占比升到 4.8\%，两个桶的占比分别变为 91.95\% 与 0.66\%，最大等价类不变，说明精度的绝大部分收益在启动期即已取得，不依赖长期运行的信息。图~\ref{fig:5-6} 中本章机制两条曲线的形状说明剩余的大等价类来自哪里：目标数在 6 到 57 之间的调用点是内核中真正的分发点，例如按文件系统类型、协议族或设备类别分发的间接调用，它们合法地调用多个函数，等价类等于目标数是这类调用点的固有属性，任何以调用点为单位的机制都不能再把它压小。与类型上下文相比，本章机制去掉的是“类型相同但从未被该调用点调用过”的那部分目标，这正是类型上下文把 141 个 \lstinline[breaklines=false]@void (*)(void)@ 函数放进同一类的原因。

\subsection{高扇入函数识别的收敛与准确性}

本节回答运行时识别需要多久收敛，以及收敛之后漏掉了多少。高扇入是一个运行时概念，静态分析在 C 语言上无法确定函数指针的实际目标集合，只能给出保守上界。

本节实验了不同时间点与负载情况下，已识别的高扇入函数数随时间的曲线、稳态数量、以及其中在运行期因绑定后归属冲突而经仲裁判为高扇入的数量。结果见表~\ref{tab:5-4} 与图~\ref{fig:5-7}，图中实线为仅靠启动期自然执行的识别数量，虚线为加入辅助负载后的数量，竖点线标出启动完成时刻。

\begin{table}[!htbp]
    \bicaption{\enspace 高扇入函数识别的收敛}{\enspace Convergence of High-Fan-In Function Identification}
    \label{tab:5-4}
    \centering
    \scriptsize
    \setlength{\tabcolsep}{3pt}
    \renewcommand{\arraystretch}{1.25}
    \hyphenpenalty=10000\exhyphenpenalty=10000\emergencystretch=1em
    \begin{tabular}{>{\raggedright\arraybackslash}p{\dimexpr 0.2134\textwidth-2\tabcolsep\relax}>{\raggedleft\arraybackslash}p{\dimexpr 0.1642\textwidth-2\tabcolsep\relax}>{\raggedleft\arraybackslash}p{\dimexpr 0.2134\textwidth-2\tabcolsep\relax}>{\raggedleft\arraybackslash}p{\dimexpr 0.1806\textwidth-2\tabcolsep\relax}>{\raggedleft\arraybackslash}p{\dimexpr 0.2134\textwidth-2\tabcolsep\relax}}
        \hline
        时刻 & 已识别高扇入函数数 & 占间接调用目标函数的比例 & 累计移入随机化区域数 & 绑定后经仲裁判为高扇入数\\
        \hline
        早期初始化 & 0 & 0\% & 0 & 0\\
        启动完成 & 233 & 4.4\% & 233 & 0\\
        运行期（UnixBench 之后） & 244 & 4.8\% & 244 & 3\\
        启动完成（辅助负载） & 261 & 5.2\% & 261 & 0\\
        运行期（辅助负载） & 264 & 5.3\% & 264 & 1\\
        \hline
    \end{tabular}
\end{table}

\begin{figure}[!htbp]
    \centering
    \includegraphics[width=0.82\textwidth]{plots/ch5-fanin-convergence}
    \bicaption{\enspace 高扇入函数识别数量随启动进程的变化}{\enspace Number of High-Fan-In Functions Identified over the Course of Boot}
    \label{fig:5-7}
\end{figure}

识别在启动期即基本收敛。图~\ref{fig:5-7} 的曲线单调递增且斜率递减：启动前 600 ms 识别出 128 个，占启动完成时的一半以上，此后只剩执行频率低的路径缓慢补入，启动完成时为 233 个。运行期一次完整的 UnixBench 只再识别出 11 个，其中 3 个是在绑定之后才被第二个调用点触及、经违规处理路径仲裁后判为高扇入并迁移的函数，说明冷路径上的高扇入函数确实存在但数量很少。静态分析给出的候选上界远大于此，因为它把所有类型匹配的目标都计入，而运行时识别只把确实被两个以上调用点触及的函数判为高扇入，两者之差正是本章不采用静态分析的原因：若按静态上界迁入，随机化区域会膨胀数倍。

未在启动期识别出的函数并非难以识别，而是启动过程中根本没有被执行到，例如磁盘配额默认未启用时的配额分配路径、启动后不再热插拔处理器时的热插拔回调、正常路径不触发的块设备请求重排队、未配置虚拟局域网时的收包分支。据此在启动期加入一组辅助负载，其作用是让上述路径上的真实调用点被执行到；表~\ref{tab:5-4} 的后两行给出其效果，启动完成时识别数由 233 增至 261，运行期绑定后的归属冲突由 3 个降到 1 个。辅助负载不能直接调用目标函数：高扇入取决于有多少个不同调用点触及该函数，若由辅助负载直接调用，只会把辅助负载自身登记为归属调用点，既识别不出扇入，还可能把一个本属单一归属的函数误判为高扇入。辅助负载的运行时机取在初始化调用序列之后、第一个用户态进程之前，使绑定仍落在安全启动链假设成立的窗口之内；其自身的调用点排除在归属统计之外；且仅在识别阶段编入，部署配置不予保留，因为它会主动触及大量平时不执行的路径。误识别率为零是机制的结构性性质：归属冲突是判定高扇入的充分条件，一个函数只有在确实被第二个调用点触及时才会被判定。

\subsection{其他部分随机化方法迁移}

本节支撑本章机制与具体随机化方法解耦这一声明。做法是把本章机制施加到一个与本文实现无关的内核部分随机化方法之上，为它留下的非随机化区域提供保护，在该方法自身的负载、间隔与资源口径下逐项对照施加前后的差别。本章机制对随机化原语的唯一要求是高扇入函数持续处于随机化状态，因此施加的步骤为：在该方法的跨区域出入口施加分区验证，把运行时识别出的高扇入函数纳入其随机化区域。

对照方法取 Chaos\citep{chang2024chaos}，理由在平台与粒度两条。本章机制依赖 ARM 指针认证，端到端测量只能在 AArch64 上进行，而 Adelie\citep{nikolaev2022adelie} 的位置无关改造建立在 x86-64 的指令指针相对寻址与全局偏移表之上，其栈重随机化还征用了该架构特有的临时寄存器，原文声明只面向 x86-64；Dbox\citep{li2024dbox} 的滑动空间与私有空间由 x86 的二级地址翻译、虚拟机扩展与输入输出虚拟化实现，原文同样声明只适用于具备这些扩展的 x86 处理器。两者都不能在本章的实验平台上运行，因此不纳入端到端测量，对它们的分析放在本节末尾。Chaos 基于 Linux 5.10 LTS 在 AArch64 上实现，以函数粒度重排内核模块，另有约 250 行的 RV64 移植，是三者中唯一能在本平台上承载本章机制的。

\textbf{两种随机化原语的对照。} 施加实验的说服力取决于对照方法与本文原型有多不同：两者越不像，解耦声明越有意义。表~\ref{tab:5-5} 按随机化原语的关键维度列出二者的差别。

\begin{table}[!htbp]
    \bicaption{\enspace 两种随机化原语的对照}{\enspace Comparison of the Two Re-randomization Primitives}
    \label{tab:5-5}
    \centering
    \scriptsize
    \setlength{\tabcolsep}{3pt}
    \renewcommand{\arraystretch}{1.25}
    \hyphenpenalty=10000\exhyphenpenalty=10000\emergencystretch=1em
    \begin{tabular}{>{\raggedright\arraybackslash}p{\dimexpr 0.1750\textwidth-2\tabcolsep\relax}>{\raggedright\arraybackslash}p{\dimexpr 0.4125\textwidth-2\tabcolsep\relax}>{\raggedright\arraybackslash}p{\dimexpr 0.4125\textwidth-2\tabcolsep\relax}}
        \hline
        维度 & Chaos\citep{chang2024chaos} & 本文原型（第 3、4 章）\\
        \hline
        随机化对象 & 内核模块，核心内核保持静态 & built-in 代码的选定子集\\
        随机化粒度 & 函数 & 函数\\
        新变体的产生 & 按函数节重排并按重定位信息重新生成整个模块 & 复制函数体到预分配的变体池中的新位置\\
        代码索引 & 每个间接地址一个固定位置的跳板 & 入口与出口跳板加跳转桩，当前地址只存在于分支立即数中\\
        执行流追踪 & 跳板前后钩子维护每模块线程表 & 跳板进出计数，区域内部的调用也计入\\
        暂停方式 & 置线程暂停标志，调度器跳过模块相关线程 & 置阻断标志，在四个入站点拦截并等待计数归零\\
        返回地址 & 每线程返回栈暂存返回地址，避免栈回溯 & 位置无关返回标记，经重入桩解析\\
        触发策略 & 固定间隔，原文取 50 ms & 固定间隔（第 3 章）与代码读取触发（第 4 章）\\
        编译期支持 & GCC 插件与跳板宏 & LLVM 插件\\
        \hline
    \end{tabular}
\end{table}

两者在新变体的产生方式、执行流追踪的口径、暂停方式与返回地址的处理这四项上都不同，唯一的共同点是随机化的粒度为函数。本章机制取用的正是这一条，它不关心变体如何产生、执行流如何暂停，只要求高扇入函数的当前地址不可预测。此外表~\ref{tab:5-5} 还暴露出一处 Chaos 自己声明为残余风险的地方：其跳板位置固定，原文承认可被用作攻击目标，理由是数量少且多由内核使用，而本章的分区验证恰好作用在这一处，施加后跳板不再是无约束的入口。

\textbf{施加前后的效果。} 表~\ref{tab:5-6} 给出施加本章机制前后的精度、开销与资源变化，端到端开销一列取 LMBench 几何均值。

\begin{table}[!htbp]
    \bicaption{\enspace 施加本章机制前后的效果}{\enspace Effect of Applying the Proposed Mechanism}
    \label{tab:5-6}
    \centering
    \scriptsize
    \setlength{\tabcolsep}{3pt}
    \renewcommand{\arraystretch}{1.25}
    \hyphenpenalty=10000\exhyphenpenalty=10000\emergencystretch=1em
    \begin{tabular}{>{\raggedright\arraybackslash}p{\dimexpr 0.2100\textwidth-2\tabcolsep\relax}>{\raggedright\arraybackslash}p{\dimexpr 0.2400\textwidth-2\tabcolsep\relax}>{\raggedleft\arraybackslash}p{\dimexpr 0.1500\textwidth-2\tabcolsep\relax}>{\raggedleft\arraybackslash}p{\dimexpr 0.1350\textwidth-2\tabcolsep\relax}>{\raggedleft\arraybackslash}p{\dimexpr 0.1350\textwidth-2\tabcolsep\relax}>{\raggedright\arraybackslash}p{\dimexpr 0.1300\textwidth-2\tabcolsep\relax}}
        \hline
        项 & 口径 & Chaos 施加前 & Chaos 施加后 & 本文原型 施加前 & 本文原型 施加后\\
        \hline
        新增随机化函数 & 个 & 0 & 233 & 0 & 233\\
        最大等价类 & 个合法目标 & 141 & 57 & 141 & 57\\
        目标数不超过 5 的调用点 & 占比 & 55.0\% & 91.0\% & 55.0\% & 91.0\%\\
        端到端开销 & LMBench 几何均值 & 2.4\% & 3.1\% & 1.98\% & 2.6\%\\
        新增常驻内存 & 迁入函数的变体与跳板 & 0 & 0.46 MB & 0 & 0.46 MB\\
        单次随机化的停顿 & 中位 & 32.7 μs & 33.1 μs & 1.22 μs & 1.24 μs\\
        \hline
    \end{tabular}
\end{table}

三项结论。其一，精度的改善与承载它的随机化方法无关：两种方法施加后的最大等价类都是 57，不超过 5 的调用点占比都是 91.0\%，因为迁入的是同一批 233 个高扇入函数，而高扇入是内核调用图的属性，不是随机化方法的属性。其二，开销增量的量级相同，为 0.7\% 与 0.6\%，来源是每次跨区域离开的 9.0 ns 验证与这 233 个函数新增的跳板经过。其三，单次随机化的停顿几乎不变，Chaos 由 32.7 μs 增至 33.1 μs、本文原型由 1.22 μs 增至 1.24 μs，因为迁入的函数只是让新变体多含 233 个函数体，而两种方法的变体都在关键路径之外生成，停顿一段只做切换。两者停顿时间的绝对值相差一个数量级，这一差别来自暂停方式本身，Chaos 需要等调度器让模块相关线程全部离开处理器，本文原型只等跳板计数归零，与本章机制无关。

\textbf{Chaos 原文负载上的对照。} 上一组数据取的是本文其余各章通用的负载，为避免口径偏向本文，本节改按 Chaos 原文的实验设置重测。负载取其所用的两个半虚拟化设备驱动 virtio-net 与 virtio-scsi，指标取相对无随机化基线的吞吐下降，随机化间隔取其所用的 50、30、10 ms 三档，每项 10 次取中位，结果见图~\ref{fig:5-8}。图中每组柱为施加本章机制前后的一对，空心方块为 Chaos 原文报告的对应值，原文以文字给出 50 ms 与 10 ms 两档。

\begin{figure}[!htbp]
    \centering
    \includegraphics[width=\textwidth]{plots/ch5-chaos-workloads}
    \bicaption{\enspace Chaos 原文四组负载在三档随机化间隔下施加本章机制前后的吞吐下降}{\enspace Throughput Degradation on the Four Workloads Used in the Chaos Paper, at Three Re-randomization Intervals, before and after Applying the Proposed Mechanism}
    \label{fig:5-8}
\end{figure}

三点可读。其一，本文平台上的复现值在四组负载、两个可比档位上都略低于原文报告值，原因是原文测于 QEMU 虚拟机而本节测于裸机，半虚拟化设备驱动的路径在裸机上更短；相对关系与原文一致，说明复现是可信的。其二，施加本章机制的增量在同一负载的三档间隔上几乎不变，virtio-net TCP 三档都是 1.0\%，UDP 为 1.1\%，virtio-scsi 读为 0.8\%、写为 0.7\%。其三，增量随负载穿越模块边界的频度上升：virtio-net 的每个报文都要穿过驱动与协议栈之间的间接调用，增量最大；virtio-scsi 写路径以批量提交为主，单位数据量上的边界穿越次数最少，增量最小。这一顺序与第 3.5.3 节对跨区域调用频度与开销关系的分析一致，说明增量由内核自身的调用频度决定，而不由随机化方法决定。

第二点值得单独说明，因为它决定了本章机制在部署上的位置。Chaos 自身的开销随间隔缩短而明显上升，virtio-net TCP 由 50 ms 档的 2.5\% 升到 10 ms 档的 7.3\%，因为每次随机化都要重新生成变体并暂停模块相关线程，这部分代价与随机化的次数成正比。本章机制的代价则按跨区域转移的次数计，与随机化的频率无关，因此增量在三档上恒定。使用者可以按自己的安全需求调整随机化间隔而不必重新评估本章机制的代价；反过来，在间隔很短因而随机化本身代价很高的配置下，本章机制的相对代价更小。

\textbf{资源开销与单次随机化延迟。} 吞吐是端到端指标，它不能说明代价落在何处。Chaos 原文另给出两项资源指标：CPU 开销随同时随机化的模块数变化，以及单次随机化中暂停并迁移模块相关线程所需的时间随线程数变化。本节在同样两项上重测，结果见图~\ref{fig:5-9}。

\begin{figure}[!htbp]
    \centering
    \includegraphics[width=\textwidth]{plots/ch5-interval-sensitivity}
    \bicaption{\enspace 施加本章机制前后的 CPU 开销与单次随机化的迁移时间}{\enspace CPU Overhead and Per-Re-randomization Migration Time before and after Applying the Proposed Mechanism}
    \label{fig:5-9}
\end{figure}

两项指标给出同一个结论。CPU 开销随同时随机化的模块数近似线性上升，由 1 个模块的 0.9\% 升到 16 个模块的 12.8\%，这一项属于 Chaos 自身：模块越多，每个周期需要生成的变体与需要暂停的线程越多。施加本章机制后每一点都抬高约 0.6\%，抬高量不随模块数变化，因为迁入随机化区域的是核心内核中固定的 233 个高扇入函数，与被随机化的模块数无关。单次随机化的迁移时间由 1 个线程的 32.7 μs 升到 16 个线程的 44.5 μs，施加本章机制后各点只增加约 0.4 μs，因为迁入的函数只是让新变体多含 233 个函数体，而变体在关键路径之外生成，迁移一段只做切换与状态重映射。两条曲线在各点上的间距基本恒定，说明本章机制的代价既不随随机化的规模放大，也不随并发度放大。

内存一项同样如此。Chaos 原文报告超过九成的内核模块的常驻内存开销不超过 240 KB，本章机制在此之上增加 0.46 MB，用于 233 个函数的两份变体与跳板，该值由迁入函数的数量决定，与模块数无关，见表~\ref{tab:5-6}。

\textbf{移植改造的规模。} 施加所需的改造量取决于两件事：该方法的随机化粒度是否已经支持函数粒度，以及它是否已有跨区域出入口可供施加验证，见表~\ref{tab:5-7}。

\begin{table}[!htbp]
    \bicaption{\enspace 向各方法移植所需的改造}{\enspace Engineering Effort Required to Port the Mechanism onto Each System}
    \label{tab:5-7}
    \centering
    \footnotesize
    \setlength{\tabcolsep}{3pt}
    \renewcommand{\arraystretch}{1.25}
    \begin{tabular}{>{\raggedright\arraybackslash}p{\dimexpr 0.1500\textwidth-2\tabcolsep\relax}>{\raggedright\arraybackslash}p{\dimexpr 0.1700\textwidth-2\tabcolsep\relax}>{\raggedright\arraybackslash}p{\dimexpr 0.3400\textwidth-2\tabcolsep\relax}r>{\raggedright\arraybackslash}p{\dimexpr 0.1700\textwidth-2\tabcolsep\relax}}
        \hline
        方法 & 随机化粒度 & 需补入的组件 & 新增代码 & 需改动的原有组件\\
        \hline
        Chaos & 模块内函数 & 分区验证 & 620 行 & 随机化名单\\
        本文原型 & built-in 函数 & 分区验证 & 540 行 & 无\\
        \hline
    \end{tabular}
\end{table}

Chaos 与本文原型只需补入分区验证，因为两者本身即函数粒度且已有跨区域出入口，新增代码都在数百行量级，其中六成以上是编译期插桩，运行期部分不超过 300 行。对 Chaos 而言，其跳板本就是所有进出模块的控制流的必经之处，分区验证直接施加在跳板的前后钩子上，不触及其变体生成与线程暂停这两条主要路径，这也是新增代码集中在编译期的原因。

\textbf{对另外两种方法的分析。} Adelie 与 Dbox 以模块或驱动为单位随机化，核心内核中的高扇入函数在其随机化区域内无处安放，施加本章机制时需要先补入第 3 章的跳板与函数粒度位置无关改造，使这 233 个函数能够以函数粒度被搬移，再在其模块边界处施加分区验证。这一路径在设计上成立，本章不给出它的测量结果，原因是两者都只能运行在 x86-64 上，而本章机制依赖 ARM 指针认证。两者与 Chaos 在本章机制看来的差别只在随机化粒度，而粒度不足这一点可由第 3 章的跳板补齐，因此这一限制来自对照方法的平台，不来自本章机制。

\section{安全性分析}

\textbf{威胁模型与假设。} 攻击者具备任意内核地址读写能力，目标是在非随机化区域内构造代码复用链，或劫持跨区域的间接转移。系统已部署一个函数粒度的持续随机化机制，本章不依赖其内部细节，只要求高扇入函数持续处于随机化状态。启动阶段不存在控制流劫持，由安全启动链保证；攻击者在劫持成功之前不能任意改写内存权限。论证以攻击者无法获得指针认证的验证预言机、且系统中不存在攻击者可控上下文的签名原语为前提。

\textbf{本章的安全保证。} 在上述前提下本章保证两条性质。其一，跨区域的间接转移其目标属于该调用点的合法目标集合：出口处攻击者需要同时伪造一个能通过调用点唯一上下文验证的认证码，且目标落在白名单内，前者的成功概率由认证码的位宽决定，每次失败的尝试都触发异常；入口处以调用形式进入随机化区域的控制流只能落在函数入口。其二，非随机化区域内部的间接转移受调用点唯一上下文约束：绑定后每个调用点只接受由它自己签发的指针，等价类等于该调用点的合法目标数，91\% 的调用点不超过 5，最大 57；高扇入函数由随机化承担，其地址不可预测。两条性质的前提有二：所有跨区域控制流转移必经出入口，该性质由随机化原语的编译器全覆盖重定向与静态扫描共同保证；归属记录在绑定后保持只读，因为绑定后归属冲突的仲裁以记录中的目标函数为依据，攻击者把指针改写为记录之外的函数时函数指针与记录不符而被拒绝。

\textbf{保证的条件性。} 推测执行侧信道可以把认证指令置于推测路径上并经微架构侧信道观测验证是否成功，从而在不产生任何崩溃、不留下任何日志的前提下爆破认证码，PACMAN\citep{ravichandran2022pacman} 已在商用处理器上跨特权级演示了这一点。本章的论证以攻击者无法获得这样的预言机为前提，针对推测执行侧信道的防御与本章正交。同类的推测执行攻击也已用于击穿 ARM 的内存标记扩展\citep{kim2024tiktag}，说明以微架构侧信道爆破硬件标签这一手法并不局限于指针认证。本章只处理前向边，返回地址的保护由第 3 章的返回标记承担，其更一般的对应物是影子栈\citep{burow2019shadow}。

\textbf{不在本章范围内的部分。} 以整函数为单位的复用：攻击者若能把某个调用点的指针改写为另一个入口跳板，仍可调用该跳板对应的整个函数，对可以进入哪些函数的约束来自调用点上下文，其强度即等价类的大小，本章不声称阻断随机化区域内以整函数为单位的复用链。返回标记的同栈深重放：以栈槽地址为上下文的签名不能阻止攻击者在构造出相同栈深度后重放同一标记，消除该残余需要链式构造。运行时识别只覆盖实际执行到的路径，冷路径上的高扇入函数在被触及之前停留在类型上下文上，其保护强度为类型级。

\section{本章小结}

部分随机化必然留下非随机化区域，其布局在运行期不变，这是所有内核持续随机化方法共同面对的问题。硬件辅助控制流完整性中精度最高的 ARM 指针认证受上下文传播限制，多个调用点调用同一函数时必须共享上下文，该约束沿调用图扩散，使实际精度大幅退化；现有工作或改进上下文构造方式，或扩展保护对象，都未能打断传播链，因为传播的根源是调用图结构本身。

本章的关键观察是随机化区域内的目标地址不可预测，进入该区域的间接调用不需要指针认证。据此本章提出以随机化打断上下文传播链的机制：以类型上下文恒定签名配合归属跟踪在运行时识别高扇入函数，不需要全局指针分析；把识别出的高扇入函数移入随机化区域，使其从传播链中被剪除，两侧调用点得以各自建立基于调用点地址的唯一上下文；在跨区域的出入口施加分区验证，进入随机化区域只验范围，离开时强制认证与白名单检查。本章对随机化原语只提出一项要求，即高扇入函数持续处于随机化状态，因而可以施加到任何一种部分随机化方法留下的非随机化区域上。

实验表明，本章机制把等价类由类型上下文下的最大 141 降到 57，91\% 的间接调用点合法目标数不超过 5，超过 50 的由 10.2\% 降到 0.76\%，代价是把 233 个高扇入函数（占间接调用目标的 4.4\%）移入随机化区域；识别在启动期即基本收敛，运行期只再增加 11 个。一次跨区域离开的开销为 9.0 ns。施加到 Chaos 之上，新增的随机化函数同样是这批 233 个，开销增量为 0.7\%，最大等价类同样降到 57。安全性论证以攻击者无法获得验证预言机为前提，推测执行侧信道类攻击不被本机制阻止；运行时识别只覆盖实际执行到的路径，冷路径上的高扇入函数在被触及之前停留在类型级保护上。

